% Part: first-order-logic
% Chapter: completeness
% Section: maximally-consistent-sets

% Definition of maximally consistent sets. Properties of mcs required
% for completeness proved are in maximally-consistent-sets.tex in the
% chapter on the proof system used.

\documentclass[../../../include/open-logic-section]{subfiles}

\begin{document}

\olfileid[hi]{fol}{com}{mcs}
\olsection{\usetoken{P}{sentence} के अधिकतम संगत समुच्चय}

\begin{defn}[अधिकतम संगत समुच्चय]
\ollabel{mcs}
!!{sentence}s का समुच्चय~$\Gamma$ \emph{अधिकतम संगत} तभी और केवल तभी है जब
\begin{enumerate}
\item $\Gamma$ संगत है, और
\item यदि $\Gamma \subsetneq \Gamma'$, तो $\Gamma'$ असंगत है।
\end{enumerate}
\end{defn}

\begin{explain}
ऊपर वाली परिभाषा के समतुल्य एक वैकल्पिक परिभाषा यह है: वाक्यों का समुच्चय
$\Gamma$ \emph{अधिकतम संगत} तभी और केवल तभी है जब
\begin{enumerate}
\item $\Gamma$ संगत है, और
\item यदि $\Gamma \cup \{ !A \}$ संगत है, तो $!A \in \Gamma$।
\end{enumerate}
दूसरे शब्दों में, किसी अधिकतम संगत समुच्चय~$\Gamma$ में पहले से न आने वाले
!!{sentence}s को~$\Gamma$ में इस तरह जोड़ना संभव नहीं है कि परिणामी
बड़ा समुच्चय संगत बना रहे।
\end{explain}

\begin{explain}
पूर्णता के प्रमाण में अधिकतम संगत समुच्चय महत्त्वपूर्ण हैं, क्योंकि हम
गारंटी दे सकते हैं कि !!{sentence}s का प्रत्येक संगत समुच्चय~$\Gamma$ किसी
अधिकतम संगत समुच्चय~$\Gamma^*$ में समाहित है, और अधिकतम संगत समुच्चय में
प्रत्येक !!{sentence}~$!A$ के लिए या तो $!A$ या उसका निषेध $\lnot !A$ होता
है। विशेषतः यह परमाण्विक !!{sentence}s के लिए सत्य है; अतः ऐसी भाषा में
जिसे !!{constant}s से उपयुक्त रूप से विस्तारित किया गया हो, किसी अधिकतम
संगत समुच्चय से हम ऐसी !!a{structure} बना सकते हैं जिसमें !!{predicate}s
की व्याख्या इस आधार पर परिभाषित होती है कि कौन-से परमाण्विक !!{sentence}s
समुच्चय~$\Gamma^*$ में हैं। फिर दिखाया जा सकता है कि यह !!{structure},
~$\Gamma^*$ के सभी !!{sentence}s को (और इसलिए $\Gamma$ में आने वाले सभी को
भी) सत्य बनाती है। इस अंतिम तथ्य के प्रमाण के लिए यह आवश्यक है कि
$\lnot !A \in
\Gamma^*$ तभी और केवल तभी जब $!A \notin \Gamma^*$, $(!A \lor !B) \in \Gamma^*$
तभी और केवल तभी जब $!A
\in \Gamma^*$ या $!B \in \Gamma^*$, इत्यादि।
\end{explain}

\begin{prop}
\ollabel{prop:mcs}
मान लें कि $\Gamma$ अधिकतम संगत है। तब:
\begin{enumerate}
\item \ollabel{prop:mcs-prov-in} यदि $\Gamma \Proves !A$, तो $!A \in
  \Gamma$।

\item \ollabel{prop:mcs-either-or} किसी भी $!A$ के लिए या तो $!A \in
  \Gamma$ या $\lnot !A \in \Gamma$।

\tagitem{prvAnd}{\ollabel{prop:mcs-and} $(!A \land !B) \in \Gamma$
  तभी और केवल तभी जब $!A \in \Gamma$ और $!B \in \Gamma$ दोनों हों।}{}

\tagitem{prvOr}{\ollabel{prop:mcs-or} $(!A \lor !B) \in \Gamma$ तभी और
  केवल तभी जब या तो $!A \in \Gamma$ या $!B \in \Gamma$ हो।}{}

\tagitem{prvIf}{\ollabel{prop:mcs-if} $(!A \lif !B) \in \Gamma$ तभी और
  केवल तभी जब या तो $!A \notin \Gamma$ या $!B \in \Gamma$ हो।}{}
\end{enumerate}
\end{prop}

\begin{proof}
निम्न सभी बातों के लिए मान लें कि $\Gamma$ अधिकतम संगत है।
\begin{enumerate}
\item यदि $\Gamma \Proves !A$, तो $!A \in \Gamma$।

मान लें कि $\Gamma \Proves !A$। इसके विपरीत मान लें कि $!A
\notin \Gamma$: तब $\Gamma$ के अधिकतम संगत होने से $\Gamma
\cup \{!A\}$ असंगत है, अतः $\Gamma \cup \{!A\} \Proves
\lfalse$।
\cref{fol:seq:prv:prop:provability-contr,fol:ntd:prv:prop:provability-contr} से
$\Gamma$ असंगत है। यह इस धारणा का खंडन करता है कि $\Gamma$ संगत है। अतः
$!A \notin
\Gamma$ नहीं हो सकता, इसलिए $!A \in \Gamma$।

\item किसी भी $!A$ के लिए या तो $!A \in \Gamma$ या $\lnot !A \in \Gamma$।

इसके विपरीत मान लें कि किसी~$!A$ के लिए $!A \notin \Gamma$ और
$\lnot !A \notin \Gamma$ दोनों हैं। चूँकि $\Gamma$ अधिकतम संगत है,
$\Gamma \cup \{!A\}$ और $\Gamma \cup \{\lnot !A\}$ दोनों असंगत हैं, अतः
$\Gamma \cup \{!A\} \Proves \lfalse$ और $\Gamma \cup
\{\lnot !A\} \Proves \lfalse$।
\cref{fol:seq:prv:prop:provability-exhaustive,fol:ntd:prv:prop:provability-exhaustive} से
$\Gamma$ असंगत है, जो विरोधाभास है। अतः ऐसा कोई !!a{sentence}~$!A$ नहीं हो
सकता और प्रत्येक $!A$ के लिए $!A \in \Gamma$ या $\lnot !A
\in \Gamma$।

\tagitem{defAnd}{}{%
\iftag{probAnd}{अभ्यास।}{%
$(!A \land !B) \in \Gamma$ तभी और केवल तभी जब $!A \in \Gamma$ और $!B \in \Gamma$
दोनों हों:

अग्र दिशा के लिए मान लें कि $(!A \land !B) \in \Gamma$। तब
$\Gamma \Proves !A \land !B$।
\cref{fol:seq:ppr:prop:provability-land-left,fol:ntd:ppr:prop:provability-land-left} से
$\Gamma \Proves !A$ और $\Gamma \Proves !B$।
\olref{prop:mcs-prov-in} से $!A \in \Gamma$ और $!B \in \Gamma$, जैसा
अपेक्षित था।

विपरीत दिशा के लिए $!A \in \Gamma$ और $!B \in
\Gamma$ मान लें। तब $\Gamma \Proves !A$ और $\Gamma \Proves !B$।
\cref{fol:seq:ppr:prop:provability-land-right,fol:ntd:ppr:prop:provability-land-right} से
$\Gamma \Proves !A \land !B$। \olref{prop:mcs-prov-in} से $(!A \land
!B) \in \Gamma$।}}

\tagitem{defOr}{}{%
\iftag{probOr}{अभ्यास।}{%
$(!A \lor !B) \in \Gamma$ तभी और केवल तभी जब या तो $!A \in \Gamma$ या
$!B \in \Gamma$ हो।

अग्र दिशा के प्रतिपक्ष के लिए मान लें कि $!A
\notin \Gamma$ और $!B \notin \Gamma$। हमें दिखाना है कि $(!A \lor
!B) \notin \Gamma$। चूँकि $\Gamma$ अधिकतम संगत है, $\Gamma
\cup \{!A\} \Proves \lfalse$ और $\Gamma \cup \{!B\} \Proves \lfalse$।
\cref{fol:seq:ppr:prop:provability-lor,fol:ntd:ppr:prop:provability-lor} से
$\Gamma \cup \{(!A \lor !B)\}$ असंगत है। अतः $(!A \lor !B)
\notin \Gamma$, जैसा अपेक्षित था।

विपरीत दिशा के लिए मान लें कि $!A \in \Gamma$ या $!B \in
\Gamma$। तब $\Gamma \Proves !A$ या $\Gamma \Proves !B$।
\cref{fol:seq:ppr:prop:provability-lor,fol:ntd:ppr:prop:provability-lor} से
$\Gamma \Proves !A \lor !B$। \olref{prop:mcs-prov-in} से $(!A \lor
!B) \in \Gamma$, जैसा अपेक्षित था।}}

\tagitem{defIf}{}{%
\iftag{probIf}{अभ्यास।}{%
$(!A \lif !B) \in \Gamma$ तभी और केवल तभी जब या तो $!A \notin \Gamma$ या
$!B \in \Gamma$ हो:

अग्र दिशा के लिए $(!A \lif !B) \in \Gamma$ मान लें, और इसके विपरीत मान
लें कि $!A \in \Gamma$ तथा $!B \notin \Gamma$। इन धारणाओं पर
$\Gamma \Proves !A \lif !B$ और $\Gamma \Proves !A$।
\cref{fol:seq:ppr:prop:provability-lif-left,fol:ntd:ppr:prop:provability-lif-left} से
$\Gamma \Proves !B$। किंतु तब \olref{prop:mcs-prov-in} से $!B \in
\Gamma$, जो $!B \notin \Gamma$ वाली धारणा का खंडन करता है।

विपरीत दिशा के लिए पहले उस स्थिति पर विचार करें जिसमें $!A \notin
\Gamma$। \olref{prop:mcs-either-or} से $\lnot !A \in \Gamma$ और इसलिए
$\Gamma \Proves \lnot !A$।
\cref{fol:seq:ppr:prop:provability-lif-right,fol:ntd:ppr:prop:provability-lif-right} से
$\Gamma \Proves !A \lif !B$। पुनः \olref{prop:mcs-prov-in} से
$(!A \lif !B) \in \Gamma$ मिलता है, जैसा अपेक्षित था।

अब उस स्थिति पर विचार करें जिसमें $!B \in \Gamma$। तब
$\Gamma \Proves !B$ और
\cref{fol:seq:ppr:prop:provability-lif-right,fol:ntd:ppr:prop:provability-lif-right} से
$\Gamma \Proves !A \lif !B$। \olref{prop:mcs-prov-in} से $(!A \lif
!B) \in \Gamma$।}}
\end{enumerate}
\end{proof}

\begin{prob}
\olref[fol][com][mcs]{prop:mcs} का प्रमाण पूरा कीजिए।
\end{prob}

\end{document}
